\begin{abstract}
Detectors of illicit activity are trained on the past and used on the future, yet they are often evaluated as if time did not matter. We study this gap with one protocol applied to three layers: Bitcoin transactions and wallets (Elliptic++) and network flows (a sample of UGR'16). Models are fitted on early time steps, tuned on later ones, and tested once on a still later period; an ensemble of gradient-boosted and bagged trees is compared with supervised and unsupervised references, and every result carries bootstrap intervals. On transactions the ensemble matches the best published result (illicit F1 0.829, 95\% interval 0.811--0.850) but is statistically indistinguishable from a 50-tree Random Forest. On wallets, features derived from the address graph double F1 (0.357 to 0.710), and graph-derived features alone match all features; the gain survives when neighbourhoods are restricted to wallets that already exist. A published wallet baseline (F1 0.845) is not reproduced under any time-ordered split (0.25--0.39) but is approached under a random split (0.87), which we trace to per-address features repeated across rows and to strong label homophily. After a known market shutdown the transaction model's F1 falls from 0.918 to 0.031; this is a ranking failure that no threshold repairs, whereas retraining restores ranking (PR-AUC 0.06 to 0.62) but leaves decisions at a fixed threshold poor. Isolation Forest scores illicit transactions as less anomalous than licit ones (ROC-AUC 0.17). The network layer is the weakest (F1 0.610) and is dominated by low-volume botnet traffic. We conclude that the split protocol, not the model family, determines the reported numbers, and that ranking and threshold failures must be reported separately.
\end{abstract}

\begin{IEEEkeywords}
anomaly detection, anti-money laundering, Bitcoin, concept drift, data leakage, graph features, network intrusion detection, reproducibility, temporal evaluation
\end{IEEEkeywords}
